\section{Differentiable phase in the reconstruction group}\label{sec:phase}

The invariant graph controls the distance to the graph above the
evolving coefficient. To select one exact trajectory on it, we
retain the reconstruction group. The following phase argument is
formulated for a finite representation coupled to an exponentially
stable equation. It is a concrete version of the asymptotic-phase
mechanism; its group-valued estimates are what will be used for
relative HYM metrics.

\subsection{The finite extension and its derivatives}
Let $V$ be a finite-dimensional real inner-product space carrying
a representation of a closed matrix group $G$. Give $G$ a
complete right-invariant Riemannian metric. A tangent variation
$\delta g$ is measured by $\xi=\delta g\,g^{-1}$ in a fixed
Lie algebra norm. Let $\beta:V\to\g$ be bounded and $C^2$, and
put $f(x)=C_x\beta(x)$. Suppose $f$ is globally Lipschitz and
the right-trivialized linearization
\begin{equation}\label{eq:finite-linearization}
 \mathcal A(x)(w,\xi)
       =(Df(x)w,D\beta(x)w+[\beta(x),\xi])
\end{equation}
satisfies $\|\mathcal A(x)\|\le L$ and
$\|D_x\mathcal A(x)\|\le K_0$.
The small quantity needed below is the upper bound $L$,
not a positive lower spectral bound.

The complete finite flow has the form
\[
 \Phi_s(x,g)=(\phi_s(x),B_s(x)g).
\]
Global Lipschitz growth controls $x$, and bounded $\beta$
controls both the matrix and its inverse on every finite time
interval. For a variation, $w=\delta x$ and
$\xi=\delta g\,g^{-1}$ satisfy
\[
 \dot w=Df(x)w,\qquad
 \dot\xi=D\beta(x)w+[\beta(x),\xi].
\]
Thus the right-trivialized derivative $T_s(x)$ depends on $x$
but not on $g$. The fundamental matrix equation, and then its
$x$ derivative, give for either sign of $s$
\begin{equation}\label{eq:finite-derivatives}
 \|T_s(x)\|\le e^{L|s|},\qquad
 \|D_xT_s(x)\|\le K_0|s|e^{2L|s|}.
\end{equation}
For the second inequality, variation of constants uses
$\|D\phi_s\|\le e^{L|s|}$; for negative times apply the
same argument to the reversed equation. Integration along
paths gives the corresponding Lipschitz bound for $\Phi_s$
in the product distance on $V\times G$.

\subsection{Statement of the phase estimate}
Let $E_1\subset E_\alpha\subset E_0$ be continuously embedded
Banach spaces, $0<\alpha<1$. Suppose $-A$ generates a
$C_0$ semigroup on $E_0$ whose restriction is also $C_0$
on $E_\alpha$, and
\begin{equation}\label{eq:abstract-semigroup}
 \|e^{-sA}\|_{E_\alpha\to E_\alpha}\le Me^{-\sigma s},
 \qquad
 \|e^{-sA}\|_{E_0\to E_\alpha}
       \le M(1+s^{-\alpha})e^{-\sigma s}.
\end{equation}
On $V\times B_{E_\alpha}(\rho)$ let
$h=(h_x,h_g)$ and $\F$ be $C^2$, with values in
$V\oplus\g$ and $E_0$, respectively, and uniform bounds
through order two. Assume
\begin{equation}\label{eq:perturbation-assumptions}
 \begin{gathered}
 h(x,0)=0,\quad \F(x,0)=0,\quad
 \|D_vh\|+\|D_v\F\|\le\epsilon,\\
 \|D_xh(x,v)\|+\|D_x\F(x,v)\|
                      \le K_0\|v\|_{E_\alpha}.
 \end{gathered}
\end{equation}
Consider
\begin{equation}\label{eq:abstract-system}
 \dot x=f(x)+h_x(x,v),\qquad
 \dot g=(\beta(x)+h_g(x,v))g,\qquad
 \dot v=-Av+\F(x,v).
\end{equation}
These bounds are in fixed finite and Banach operator norms;
none is a bound in the absolute matrix coordinates of $g$.

\begin{proposition}\label{prop:phase}
Fix the constants in \eqref{eq:abstract-semigroup} and the
uniform $C^2$ bounds. If $L,\epsilon$ are sufficiently small,
there is $\delta>0$ such that for
$D=\|v_0\|_{E_\alpha}<\delta$, \eqref{eq:abstract-system}
has a global mild solution and a phase
\begin{equation}\label{eq:phase-limit}
 \Pi(x_0,g_0,v_0)
        =\lim_{T\to\infty}\Phi_{-T}(x(T),g(T))
          \quad\hbox{in } V\times G.
\end{equation}
The map $\Pi$ is $C^1$ with uniform derivative bounds in
right-trivialized norms. For $a=\sigma/2$ and
$Y_*(s)=\Phi_s\Pi$,
\begin{equation}\label{eq:abstract-tracking}
 \|v(s)\|_{E_\alpha}\le CD e^{-as},\qquad
 d((x(s),g(s)),Y_*(s))\le C\epsilon D e^{-as}.
\end{equation}
The selected trajectory is unique among complete finite
trajectories tracking at $O(e^{-qs})$ for some $q>L$.
If $v_0=0$, its phase is $(x_0,g_0)$.
The constants are uniform for families with the same assumptions,
independent of $g_0$ and of finite spectral gaps.

If $h_x=C_xh_g$ and $x_0=g_0x_b$, then
$x_*(s)=g_*(s)x_b$.
\end{proposition}

\begin{proof}
\emph{Normal decay and continuation.}
The weighted convolution constant
\[
 J_a=M\int_0^\infty(1+t^{-\alpha})e^{-(\sigma-a)t}\,dt
\]
is finite. Duhamel's formula implies
\[
 \sup_{s\le T}e^{as}\|v(s)\|
 \le MD+\epsilon J_a\sup_{s\le T}e^{as}\|v(s)\|.
\]
Choose $\epsilon J_a\le1/2$ and $2M\delta<\rho/2$.
Then $\|v(s)\|\le2MD e^{-as}$, so the solution never
reaches the normal boundary. The globally Lipschitz finite
field and bounded perturbations exclude finite-time escape
of $x$ or $g$. Local mild existence and continuation follow
by contraction with the integrable smoothing kernel
$t^{-\alpha}$. The same argument gives local $C^1$
dependence in the stated initial topology.

\smallskip
\emph{Convergence in the group and tracking.}
Write $Y(s)=(x(s),g(s))$ and $p_T=\Phi_{-T}Y(T)$.
Its velocity in the invariant frames at $p_T$ is
\begin{equation}\label{eq:phase-velocity}
 \dot p_T=T_{-T}(x(T))h(x(T),v(T))=:w_T.
\end{equation}
By \eqref{eq:finite-derivatives},
$\|w_T\|\le2M\epsilon D e^{-(a-L)T}$.
For $L<a$, the curve $p_T$ has finite length, so
completeness of $V\times G$ gives \eqref{eq:phase-limit}.
Its distance from $(x_0,g_0)$ is at most
$2M\epsilon D/(a-L)$.

For fixed $s$, differentiate $\Phi_{s-T}Y(T)$ for
$T\ge s$. The velocity bound is
\[
 2M\epsilon D e^{L(T-s)}e^{-aT}.
\]
Integration gives
\[
 d(Y(s),\Phi_s\Pi)
       \le \frac{2M\epsilon D}{a-L}e^{-as}.
\]
If two finite trajectories have distance $O(e^{-qs})$
with $q>L$, applying $\Phi_{-s}$ forces their initial
distance to be zero. A second trajectory tracking the
actual solution at such a rate differs from the selected
one at rate $\min(a,q)>L$, proving the stated uniqueness.
Existence at every rate $q>L$ is not asserted.

\smallskip
\emph{Differentiability of the limit.}
The $(x,v)$ system is independent of $g$. Set
$\gamma=\sigma/16$ and require $L<\gamma/4$.
For variations of its initial data write $W=\delta x$
and $Z=\delta v$. Then
\[
 \dot W=(Df+D_xh_x)W+D_vh_xZ
\]
and
\[
 Z(s)=e^{-sA}Z_0+\int_0^s e^{-(s-t)A}
       (D_x\F\,W+D_v\F\,Z)(t)\,dt.
\]
Define
$U_T=\sup_{s\le T}e^{-\gamma s}|W(s)|$ and
$V_T=\sup_{s\le T}e^{(a-\gamma)s}\|Z(s)\|$.
The normal convolution estimate, with $\epsilon$ smaller
if necessary, gives
\[
 V_T\le C\|Z_0\|+CK_0D U_T.
\]
Variation of constants for $W$ uses
$\int_0^\infty\|v(t)\|dt\le2MD/a$, so
$D_xh_x$ has integrable norm. It gives
$U_T\le C|W_0|+C\epsilon V_T$.
Choose $\delta$ so that the resulting
$C\epsilon K_0D U_T$ term is absorbed. Uniformly in $T$,
\begin{equation}\label{eq:initial-variations}
 \begin{split}
 |W(s)|&\le Ce^{\gamma s}(|W_0|+\epsilon\|Z_0\|),\\
 \|Z(s)\|&\le Ce^{-(a-\gamma)s}(\|Z_0\|+D|W_0|).
 \end{split}
\end{equation}
The convolution is finite because $a-\gamma<\sigma$.

Differentiate the right-trivialized velocity
\eqref{eq:phase-velocity} with respect to $(x_0,v_0)$:
\begin{equation}\label{eq:phase-velocity-derivative}
 \delta w_T=(D_xT_{-T})(x(T))[W(T)]h(x(T),v(T))
       +T_{-T}(x(T))(D_xh\,W+D_vh\,Z).
\end{equation}
Its norm is bounded by a constant times
\begin{multline}\label{eq:integrable-derivative}
 \bigl[\epsilon D(1+T)e^{-(a-2L-\gamma)T}
       +D e^{-(a-L-\gamma)T}
       +\epsilon e^{-(a-L-\gamma)T}\bigr]\\
 {}\cdot(|W_0|+\|Z_0\|).
\end{multline}
This bound is integrable, uniformly on a smaller initial
normal ball.

To pass from this estimate to derivatives of $p_T$, use
its Euclidean component and an identity chart for its
group component after right translation by $g_0^{-1}$.
The phase-length estimate keeps
$q_T=p_T^Gg_0^{-1}$ inside that chart for small $\delta$.
It satisfies $\dot q_T=w_T^Gq_T$ and
\[
 \partial_T(\delta q_Tq_T^{-1})
       =[w_T^G,\delta q_Tq_T^{-1}]+\delta w_T^G.
\]
Integrability of $w_T$ and $\delta w_T$ implies uniform
convergence of these derivatives. The limits are continuous
by dominated convergence and finite-time continuous
dependence. Thus the limiting phase is $C^1$ in $(x_0,v_0)$.
For arbitrary $g_0$, right equivariance gives
\[
 \Pi^G(x_0,g_0,v_0)=\Pi^G(x_0,I,v_0)g_0.
\]
Its derivative in $g_0$, in right-trivialized norms, is
conjugation by an element in the same fixed identity
chart, and is uniformly bounded.

\smallskip
\emph{Preservation of the orbit.}
When $h_x=C_xh_g$, both perturbed and unperturbed
equations preserve $x=gx_b$. Every finite-time phase
point satisfies $p_T^V=p_T^Gx_b$. Since the limit
converges in the group itself, this equality passes
to the limit and then along the finite flow.
No inverse of $C_x$ occurs.
\end{proof}
